% Generated draft: requires mathematical and source review.
\documentclass[11pt]{article}
\usepackage[margin=25mm]{geometry}
\usepackage{fontspec}
\setmainfont{Latin Modern Roman}
\usepackage{amsmath,amssymb,mathtools,unicode-math}
\setmathfont{Latin Modern Math}
\usepackage{enumitem,needspace,xurl,hyperref}
\setlength{\parindent}{0pt}
\setlength{\parskip}{5pt}
\setlength{\emergencystretch}{3em}
\newcommand{\N}{\mathbb{N}}
\newcommand{\Z}{\mathbb{Z}}
\newcommand{\Q}{\mathbb{Q}}
\newcommand{\R}{\mathbb{R}}
\newcommand{\C}{\mathbb{C}}
\newcommand{\F}{\mathbb{F}}
\newcommand{\A}{\mathbb{A}}
\newcommand{\PP}{\mathbb{P}}
\newcommand{\E}{\mathbb{E}}
\newcommand{\eps}{\varepsilon}
\newcommand{\dd}{\,\mathrm{d}}
\begin{document}
\section*{UnsolvedMath ID: 30001839}
\textbf{Source problem code:} catalogue entry
\subsection{Definitions and mathematical statement}
Let \(G\) be a fixed finite group. Let \(L/K\) be a finite Galois extension, let \(\Gamma=\operatorname{Gal}(L/K)\), and let \(\iota\colon\Gamma\hookrightarrow G\) be an injective homomorphism. Write \(H=\iota(\Gamma)\). The desired conclusion is that \(\iota\) is an isomorphism, equivalently \(H=G\).

The supplied question asks which finite collection of conjugacy-class observations or Frobenius data suffices to certify this conclusion. One precise interpretation of conjugacy-class observations is to seek a nonnegative integer \(r\) and conjugacy classes \(C_1,\ldots,C_r\) of \(G\) such that
\[
\forall J\leq G,\qquad
\left(\forall j\in\{1,\ldots,r\},\ J\cap C_j\ne\varnothing\right)
\Longrightarrow J=G.
\]
Under this interpretation, the observations are the existence of elements \(\gamma_j\in\Gamma\) satisfying \(\iota(\gamma_j)\in C_j\) for each \(j\).

This interpretation is a formulation for review, not an identification of an unambiguously specified source target. The record does not say whether observations mean intersection with prescribed classes, exact element data, class frequencies, or another type of information; nor does it specify an optimality requirement, an algorithm, or permitted dependence of the certificate on the extension. Quantifying over every subgroup \(J\leq G\) expresses a uniform group-theoretic certificate; a certificate restricted to a particular arithmetic family could have a different quantifier range.

For a Frobenius formulation, the record does not specify the arithmetic nature of \(K\), the places to be used, or the hypotheses ensuring that the relevant Frobenius conjugacy classes are defined. Accordingly, no particular field domain, unramified-place condition, or bound on places is silently imposed here. A precise arithmetic version remains unresolved from the supplied record.
\subsection{Short English statement}
Given an embedding of a finite Galois group into a finite group \(G\), determine what finite conjugacy-class or Frobenius observations certify that its image is all of \(G\). The displayed condition formalizes one possible interpretation using classes met by the image; the intended observation model and arithmetic hypotheses require clarification.
\subsection{Sources}
\begin{itemize}
\item Supplied local UnsolvedMath registry record, permanent ID \(30001839\), title \emph{Frobenius Certificates for Finite Galois-Group Isomorphisms}, source problem code \texttt{catalogue entry}, in the Oberwolfach Reports Open Problems collection. The record gives proposed year 2011 and names no individual proposer.
\item Supplied report citation: \emph{Explicit Methods in Number Theory (2012)}, Oberwolfach Reports, pp.~1971--2044, DOI \texttt{10.4171/owr/2011/35}. \url{https://doi.org/10.4171/owr/2011/35}. The collection describes these reports as documenting mathematical workshops and their open-problem discussions. The supplied report title/year and DOI are retained without independent bibliographic verification.
\item Supplied literature reference: Kowalski--Zywina, ``Chebotarev invariant.'' \url{https://arxiv.org/abs/1008.4909}.
\item Supplied literature reference: Jouve--Kowalski--Zywina, ``finite-group applications.'' \url{https://doi.org/10.5802/jtnb.649}.
\item Literature assessment dated exactly \texttt{2026-08-21}: the supplied curation labels the problem partially solved and states that Chebotarev and Frobenius-conjugacy methods give effective tests in important settings, while the broad extracted formulation is not completely classified. This is supplied dated curation, not an independently verified assessment; the record itself requests independent verification of time-sensitive status and cited claims.
\item The registry labels the statement \texttt{corrected\_verified} and reports that its mathematical content was checked against the official Oberwolfach report, with notation, context, and question form normalized without changing the source-level claim. This verification claim is supplied metadata and has not been independently checked here.
\end{itemize}
\end{document}
